\chapter{Workflow-Based Baseline}
\label{ch:workflow}

% The control condition. It must be COMPETENT, not a strawman: a faithful
% implementation of an established coordination doctrine, frozen before the
% novel disruptions are revealed. See CONTEXT.md, "Workflow Baseline".

\section{Design Rationale}
\label{sec:wf:rationale}
% Why a plain Python script and not Camunda/Temporal/Airflow: the engine is not
% the independent variable. Pre-scripted branching is, see
% Section~\ref{sec:bg:workflow}.

\section{Doctrine Grounding}
\label{sec:wf:doctrine}
% DECIDE FIRST: DV 100 or ICS as the single primary anchor. CONTEXT.md still
% lists both as candidates.
% Element-by-element mapping table: doctrine element -> baseline component, each
% row cited. Design criterion is fidelity to the source model, explicitly NOT
% "fit to our storyline". Combination allowed only where the primary model lacks
% a coordination step, and each borrowed element sourced.

\section{Sequential Ten-Step Workflow}
\label{sec:wf:steps}
% TODO: Walk through each of the ten storyline steps as procedural code.
% Cross-reference Section~\ref{sec:cs:storyline}.

\section{Branching for Anticipated Disruptions}
\label{sec:wf:disruptions}
% The pre-written handlers. Precondition gate: the baseline must pass 100% of
% the anticipated disruption set before it is allowed to freeze.

\section{Freeze Protocol}
\label{sec:wf:freeze}
% The audit trail. Baseline committed and tagged in git BEFORE evaluation-day
% disruptions are revealed; record commit hash and date here.
% This is the firewall evidence that the baseline could not have pre-encoded the
% novel disruptions. Optional co-certification by the domain expert.
% Artifact recorded in Appendix~\ref{app:frozen}.

\section{Scenario Runner}
\label{sec:wf:cli}
% How a run is launched, and how it reads the same scenario inputs and writes to
% the same event log as the agentic system.

\section{Expected Failure Mode}
\label{sec:wf:limits}
% NOT a list of weaknesses built in on purpose. State it as a principled property
% of pre-scripted branching: anticipated exceptions get build-time handlers,
% unanticipated ones fall back to a generic or degraded response
% (Reichert & Weber). Note the compensating strengths the baseline is expected to
% win on: latency, cost, and determinism.
